\subsection{Agent bazowy}
\label{sec:baseline}

Planowane są dwa agenty bazowe:
\begin{itemize}
  \item \textbf{agent losowy} -- wybiera akcję z rozkładu jednostajnego;
        wyznacza dolną granicę wyników,
  \item \textbf{agent zachłanny} -- wyznacza najkrótszą ścieżkę (BFS) do
        najbliższego niesfotografowanego celu, a po zebraniu wszystkich celów
        -- do bazy; ignoruje baterię, radary i patrol. Jest to uczciwy
        punkt odniesienia, ponieważ reprezentuje „rozsądną'' strategię
        bez uczenia.
\end{itemize}

\todo{Etap E2: wyniki agentów bazowych na każdej mapie.}
